%----------------------------------------------------------------------------------------
%   Доорх хэсгийг өөрчлөх шаардлагагүй
%----------------------------------------------------------------------------------------
%!TEX TS-program = xelatex
%!TEX encoding = UTF-8 Unicode
\documentclass[12pt,a4paper]{report}

\usepackage{fontspec,xltxtra,xunicode}
\usepackage{textcomp}
\setmainfont[Ligatures=TeX]{Times New Roman}
\setsansfont{Arial}

% \usepackage[utf8x]{inputenc}
% \usepackage[mongolian]{babel}
%\usepackage{natbib}
\usepackage{geometry}
%\usepackage{fancyheadings} fancyheadings is obsolete: replaced by fancyhdr. JL
\usepackage{fancyhdr}
\setlength{\headheight}{15pt}
\usepackage{float}
\usepackage{afterpage}
\usepackage{graphicx}
\usepackage{amsmath,amssymb,amsbsy}
\usepackage{dcolumn,array}
\usepackage{longtable}
\usepackage{booktabs}
\usepackage{tabularx}
\usepackage{ltablex}
\usepackage{tocloft}
\usepackage{dics}
\usepackage{microtype}
\usepackage{ragged2e}
\usepackage{nomencl}
\usepackage{upgreek}
\newcommand{\argmin}{\mathop{\mathrm{arg\,min}}}
\usepackage{mathtools}
\usepackage[hidelinks]{hyperref}
\usepackage{tikz}
\usepackage{algorithm}
\usepackage{algpseudocode}
\usepackage{xcolor}
\usepackage{listings}
\DeclarePairedDelimiter\abs{\lvert}{\rvert}%
\makeatletter
\setlength{\belowcaptionskip}{0.5pt}
\usepackage{subfiles}
\usepackage{url}
\usepackage{xurl}
\keepXColumns
% \usepackage{minted}
% \usepackage[utf8]{inputenc}
% \setminted{
%     fontsize=\small,
%     breaklines=true,
%     frame=single,
%     numbers=left,
%     style=friendly
% }
\renewcommand{\lstlistingname}{Код}
\renewcommand{\lstlistlistingname}{\lstlistingname ын жагсаалт}

\usepackage{color}
\definecolor{codegreen}{rgb}{0,0.6,0}
\definecolor{codegray}{rgb}{0.5,0.5,0.5}
\definecolor{codepurple}{rgb}{0.58,0,0.82}
\definecolor{backcolour}{rgb}{0.99,0.99,0.99}


\lstdefinestyle{mystyle}{
    basicstyle=\ttfamily\small,
    backgroundcolor=\color{backcolour},   
    commentstyle=\color{codegreen},
    keywordstyle=\color{magenta},
    numberstyle=\tiny\color{codegray},
    stringstyle=\color{codepurple},
    basicstyle=\footnotesize,
    breakatwhitespace=false,         
    breaklines=true,                 
    columns=fullflexible,
    captionpos=b,
    keepspaces=true,
    numbers=left,
    numbersep=10pt,                  
    showspaces=false,                
    showstringspaces=false,
    showtabs=false,
    tabsize=2,
    postbreak=\mbox{\textcolor{codegray}{$\hookrightarrow$}\space}
}
 
\lstset{style=mystyle}

\raggedbottom
\emergencystretch=3em
\tolerance=2000
\hbadness=2000
\hfuzz=1pt
\Urlmuskip=0mu plus 1mu

\newcolumntype{L}[1]{>{\raggedright\arraybackslash}p{#1}}
\newcolumntype{Y}{>{\RaggedRight\arraybackslash}X}
\newcommand{\codeword}[1]{\begingroup\small\ttfamily\path{#1}\endgroup}

\pagestyle{fancy}
\fancyhf{}
\fancyhead[L]{\small\nouppercase{\leftmark}}
\fancyhead[R]{}
\fancyfoot[C]{\thepage}
\renewcommand{\headrulewidth}{0pt}

\makeatletter
\renewcommand{\chaptermark}[1]{\markboth{\thechapter.\ #1}{}}
\renewcommand{\sectionmark}[1]{}
\makeatother

\fancypagestyle{plain}{%
    \fancyhf{}
    \fancyfoot[C]{\thepage}
    \renewcommand{\headrulewidth}{0pt}
}

\let\oldabs\abs
\def\abs{\@ifstar{\oldabs}{\oldabs*}}
\makenomenclature
\begin{document}


%----------------------------------------------------------------------------------------
%   Өөрийн мэдээллээ оруулах хэсэг
%----------------------------------------------------------------------------------------

% Дипломийн ажлын сэдэв
\title{Backend системд зориулсан annotation-д суурилсан Java framework-ийн загварчлал, хэрэгжүүлэлт}
% Дипломын ажлын англи нэр
\titleEng{Design and Implementation of an Annotation-Based Java Framework for Backend Systems
}
% Өөрийн овог нэрийг бүтнээр нь бичнэ
\author{Гантулгын Батням}
% Өөрийн овгийн эхний үсэг нэрээ бичнэ
\authorShort{Г.Батням}
% Удирдагчийн зэрэг цол овгийн эхний үсэг нэр
\supervisor{Др. Т.Цэрэннадмид}
% Хамтарсан удирдагчийн зэрэг цол овгийн эхний үсэг нэр
\cosupervisor{Др. Т.Цэрэннадмид}

% СиСи дугаар 
\sisiId{22B1NUM5578}
% Их сургуулийн нэр
\university{МОНГОЛ УЛСЫН ИХ СУРГУУЛЬ}
% Бүрэлдэхүүн сургуулийн нэр
\faculty{МЭДЭЭЛЛИЙН ТЕХНОЛОГИ, ЭЛЕКТРОНИКИЙН СУРГУУЛЬ}
% Тэнхимийн нэр
\department{МЭДЭЭЛЭЛ, КОМПЬЮТЕРЫН УХААНЫ ТЭНХИМ}
% Зэргийн нэр
\degreeName{Бакалаврын судалгааны ажил}
% Суралцаж буй хөтөлбөрийн нэр
\programeName{Программ хангамж(D 061302)}
% Хэвлэгдсэн газар
\cityName{Улаанбаатар}
% Хэвлэгдсэн огноо
\gradyear{2026 оны 4 дугаар сар}


%----------------------------------------------------------------------------------------
%   Доорх хэсгийг өөрчлөх шаардлагагүй
%----------------------------------------------------------------------------------------
%----------------------Нүүр хуудастай хамаатай зүйлс----------------------------
\pagenumbering{roman}
\makefrontpage
\maketitle

\doublespace

\markboth{}{}
\pagestyle{plain}

% Decleration
\begin{huge}
\textbf{Зохиогчийн баталгаа}
\end{huge}
\par\vspace{\baselineskip}
\doublespace
\noindent
Миний бие \thesisAuthor\ ``\textbf{\thesisTitle}'' сэдэвтэй
\textbf{\thesisDegreeName}-ыг гүйцэтгэсэн болохыг зарлаж,
дараах зүйлсийг баталж байна:
\begin{itemize}
\item Ажил нь бүхэлдээ Монгол Улсын Их Сургуульд дээд боловсролын зэрэг горилохоор дэвшүүлсэн болно.
\item Энэ ажлын аль нэг хэсгийг эсвэл бүхлээр нь ямар нэг их, дээд сургуулийн зэрэг горилохоор оруулж байгаагүй болно.
\item Бусдын хийсэн ажлаас хуулбарлаагүй, эшлэл, зүүлтийг зохистой хийсэн болно.
\item Ажлыг зохиогч би хийсэн ба миний хийсэн ажил, бусдын дэмжлэгийг дипломын ажилд тодорхой тусгасан болно. 
\item Ажилд тусалсан бүх эх сурвалжид талархаж байна. 
\end{itemize} 
\par

Гарын үсэг: \underline{\hspace{5cm}} 

Огноо: 	\ \ \underline{\hspace{3cm}}

% Гарчгийг автоматаар оруулна
\setcounter{tocdepth}{1}
\tableofcontents

% Зургийн жагсаалтыг автоматаар оруулна
\listoffigures

% Хүснэгтийн жагсаалтыг автоматаар оруулна
\listoftables

% Кодын жагсаалтыг автоматаар оруулна
\lstlistoflistings

\clearpage
\markboth{}{}
\pagestyle{plain}

% This puts the word "Page" right justified above everything else.
%% \addtocontents{lof}{Зураг~\hfill Хуудас \par}
%% \addtocontents{lot}{Хүснэгт~\hfill Хуудас \par}

\renewcommand{\cftlabel}{Зураг}


\doublespace
\pagenumbering{arabic}


\clearpage
\pagestyle{plain}
\markboth{}{}
\markright{}

% Удиртгалыг оруулж ирэх ба intro.tex файлд удиртгалаа бичнэ
\cleardoublepage
\thispagestyle{plain}
\addcontentsline{toc}{chapter}{Удиртгал}
\begin{center}
    \vspace*{-1.2\baselineskip}
    {\huge\bfseries Удиртгал\par}
\end{center}
\vspace{0.5\baselineskip}
\subfile{intro}

%----------------------------------------------------------------------------------------
%   Дипломын үндсэн хэсэг эндээс эхэлнэ
%----------------------------------------------------------------------------------------
%\addcontentsline{toc}{part}{БҮЛГҮҮД}
\clearpage
\markboth{}{}
\markright{}
\pagestyle{fancy}
% Шинэ бүлэг
\subfile{research.tex}
\subfile{requirement.tex}
\subfile{design_architecture.tex}
\subfile{implementation}
\subfile{summer.tex}

\singlespace
\clearpage
\markboth{НОМ ЗҮЙ}{}
\addcontentsline{toc}{part}{НОМ ЗҮЙ}
\begin{thebibliography}{}

    \bibitem{image1}
    Inserting Images, Share LaTex,
    \url{https://www.sharelatex.com/learn/Inserting_Images}

    \bibitem{pharagraph1}
    Paragraphs and new lines, Share LaTex,
    \url{https://www.sharelatex.com/learn/Paragraphs_and_new_lines}

    \bibitem{format1}
    Bold, italics and underlining, Share LaTex,
    \url{https://www.sharelatex.com/learn/Bold,_italics_and_underlining}

    \bibitem{list}
    Lists, Share LaTex,
    \url{https://www.sharelatex.com/learn/Lists}

    \bibitem{table}
    Tables, Share LaTex,
    \url{https://www.sharelatex.com/learn/Tables}

    \bibitem{fowler2004injection}
    Fowler, M. (2004).
    \textit{Inversion of Control Containers and
    the Dependency Injection Pattern}.
    \url{https://martinfowler.com/articles/injection.html}

    \bibitem{fowler2004ioc}
    Fowler, M. (2005).
    \textit{Inversion of Control}. Martin Fowler's Bliki.
    \url{https://martinfowler.com/bliki/InversionOfControl.html}

    \bibitem{spring2024ioc}
    Spring Framework Documentation. (2024).
    \textit{The IoC Container}.
    \url{https://docs.spring.io/spring-framework/reference/core/beans.html}

    \bibitem{spring2024annotation}
    Spring Framework Documentation. (2024).
    \textit{Annotation-based Container Configuration}.
    \url{https://docs.spring.io/spring-framework/reference/core/beans/annotation-config.html}

    \bibitem{spring2024xml}
    Spring Framework Documentation. (2024).
    \textit{XML-based Container Configuration}.
    \url{https://docs.spring.io/spring-framework/reference/core/beans/basics.html}

    \bibitem{spring2024java}
    Spring Framework Documentation. (2024).
    \textit{Java-based Container Configuration}.
    \url{https://docs.spring.io/spring-framework/reference/core/beans/java.html}

    \bibitem{spring2024beanfactory}
    Spring Framework Documentation. (2024).
    \textit{The BeanFactory API}.
    \url{https://docs.spring.io/spring-framework/reference/core/beans/beanfactory.html}

    \bibitem{oracle2024annotation}
    Oracle. (2024).
    \textit{Lesson: Annotations}.
    The Java\texttrademark{} Tutorials.
    \url{https://docs.oracle.com/javase/tutorial/java/annotations/}

    \bibitem{oracle2024target}
    Oracle. (2024).
    \textit{Annotation Type Target}.
    Java SE 17 API Documentation.
    \url{https://docs.oracle.com/en/java/javase/17/docs/api/java.base/java/lang/annotation/Target.html}

    \bibitem{oracle2024declaring}
    Oracle. (2024).
    \textit{Declaring an Annotation Type}.
    The Java\texttrademark{} Tutorials.
    \url{https://docs.oracle.com/javase/tutorial/java/annotations/declaring.html}

    \bibitem{guice2024}
    Google. (2024).
    \textit{Google Guice: A lightweight dependency injection
    framework for Java}. GitHub.
    \url{https://github.com/google/guice/wiki/Motivation}

    \bibitem{jakarta2024cdi}
    Jakarta EE. (2024).
    \textit{Jakarta Contexts and Dependency Injection 4.1}.
    \url{https://jakarta.ee/specifications/cdi/4.1/}

    \bibitem{maven2024}
    Apache Software Foundation. (2024).
    \textit{Maven -- Welcome to Apache Maven}.
    \url{https://maven.apache.org/index.html}

    \bibitem{jetbrains2024}
    JetBrains. (2024).
    \textit{IntelliJ IDEA: The Leading Java and Kotlin IDE}.
    \url{https://www.jetbrains.com/idea/}

    \bibitem{jetbrains2023survey}
    JetBrains. (2023).
    \textit{The State of Developer Ecosystem 2023: Java}.
    JetBrains Developer Survey.
    \url{https://www.jetbrains.com/lp/devecosystem-2023/java/}

    \bibitem{gamma1994}
    Gamma, E., Helm, R., Johnson, R., \& Vlissides, J. (1994).
    \textit{Design Patterns: Elements of Reusable
    Object-Oriented Software}. Addison-Wesley.

    \bibitem{visualparadigm}
    Visual Paradigm Online.
    \url{https://online.visual-paradigm.com}

\end{thebibliography}

\subfile{code}
\end{document}
